% ============================================
%         Preset 1. Document Type
% ============================================
\documentclass[lettersize,journal]{IEEEtran}

% ============================================
%         Preset 2. Local Path
% ============================================
% When you create "figures" or "movies" directories in the "src" (source code) directory, the file name including path can be long.
% To avoid this, and to make it easier to adjust, you can define the alias for the path.
% The examples are provided below.

\newcommand*{\FIGURESPATH}{./figures}
% \newcommand*{\SIMFIGURESPATH}{./src/script_simulation/figures}
% \newcommand*{\SLXFIGURESPATH}{./src/simulink_simulation/figures}
% \newcommand*{\MOVIESPATH}{./movies}

% ============================================
%         Preset 3. Pre-defined Settings
% ============================================
% PLEASE DO NOT ADD OR REMOVE PACKAGES IN THE SUBMODULE LOCALLY!
% CONTACT THE AUTHOR FOR ADJUSTMENTS.
%
% The packages are pre-defined in the submodule "Template".
% If you need more packages, please add them after using pre-defined packages.

\def\cmt{true} % true for comments, false for no comments
\def\pub{false} % true for publication, false for draft
\newcommand*{\template}{../../template} 
\input{\template/preamble/preamble_journal.tex}

% ============================================
%     Preset 4. Additional Corrections
% ============================================
% correct bad hyphenation here
% \hyphenation{op-tical net-works semi-conduc-tor}
% \pagestyle{empty}

\begin{document}

% ============================================
%            TITLE and AUTHORS
% ============================================
\title{
  Vector-Space Optimization for Contraction Theory-Based Control Design: An Energy-Based Effective Space Approach
} %% Article title

\author{
  Myeongseok Ryu, Sesun You, and Kyunghwan Choi
\thanks{
  Myeongseok Ryu and Kyunghwan Choi are with the Cho Chun Shik Graduate School of Mobility, Korean Advanced Institute of Science and Technology (KAIST), Daejeon 34051, Korea (e-mail: dding\_98@kaist.ac.kr and kh.choi@kaist.ac.kr).
  Sesun You is with the \cred{Department} of Incheon National University, Incheon 22012, Korea (e-mail: yousesun@inu.ac.kr).
  }
}

\markboth{Journal of \LaTeX\ Class Files,~Vol.~18, No.~9, September~2020}%
{How to Use the IEEEtran \LaTeX \ Templates}

\maketitle

% ============================================
%         ABSTRACT and KEYWORDS
% ============================================
\begin{abstract}
  \cblue{\lipsum[1]}
\end{abstract}

\begin{IEEEkeywords}
  Contraction theory, 
  optimization-based control,
  % input saturation, 
  energy-based control
\end{IEEEkeywords}

% ============================================
%         Notation
% ============================================
\section*{Notation}
In this study, the following notation is used:

\begin{itemize}
    \item $\R^n$ denotes the $n$-dimensional Euclidean space.
    \item $\R^{n\times m}$ denotes the set of $n\times m$ real matrices.
    \item $\otimes$ denotes the Kronecker product \cite[Chap. 7 Def. 7.1.2]{Bernstein:2009aa}.
    \item A vector and a matrix are denoted by $\mv{x}=[x_i]_{i\in\{1,\cdots,n\}}\in\R^n$ and $
        \mm A
        := 
        [a_{ij}]
        _{
            i\in\{1,\cdots,n\},j\in\{1,\cdots ,m\}
        }\allowbreak\in\R^{n\times m}
        $, respectively.
    \item $\myrow_i(\mm A)$ denotes the $i\textsuperscript{th}$ row of the matrix $\mm A\in\R^{n\times m}$. 
    \item For $\mm{A}\in\R^{n\times m}$, denotes the vectorization of $\mm A$, $\myvec(\mm A):=(\myrow_1(\mm A^\top),\cdots,\myrow_m(\mm A^\top)  )^\top\in\R^{nm}$.
    % \item $\mycol_i(A)$ and $\myrow_j(A)$ denote the $i\textsuperscript{th}$ column and $j\textsuperscript{th}$ row of $A\in\R^{n\times m}$, respectively.
    % \item $\myvec(A):= [\mycol_1(A)^\top  ,\cdots,\mycol_m(A)^\top  ]^\top   $ for $A\in\R^{n\times m}$.
    \item $\lambda_{\min}(\mm A)$ denotes the minimum eigenvalue of the matrix $\mm A\in\R^{n\times n}$.
    \item $\mm I_n$ denotes the $n\times n$ identity matrix, and $\mm 0_{n\times m}$ denotes the $n\times m$ zero matrix.
\end{itemize}

% ============================================
%         SECTION: Dummy Section
% ============================================
\section{Introduction}
\label{sec:intro}

\subsection{Motivation}

\subsection{Related Works}

\subsection{Contributions}

\subsection{Organization}

\section{Preliminaries}

\section{Problem Formulation}

\section{Main Results}

\subsection{Projection of Contraction Condition}

In this section, we present the proposed vector-space optimization framework.
We project the contraction condition \eqref{eq:cstr:ctrc} onto the trajectory error vector space by pre- and post-multiplying it by $\mv{e}^\top$ and $\mv{e}$, respectively, as follows:
\begin{equation} \label{eq:cstr:ctrc:proj}
   -
   2
   \mv{y}^\top 
   \left(
      \mm{B} 
      \mm{R}^{-1} 
      \mm{B}^\top
   \right)
   \mv{y}
   +
   \mv{e}^\top
   \!
   \left(
      \tfrac{1}{T_s}+2\alpha
      +
      2\mathbb{A}^\top
   \!
   \right)
   \mv{y}
   -
   \tfrac{\mv{e}^\top \mm{M}_{\rm{pre}} \mv{e}}{T_s}
   \le 0
   ,
\end{equation}
where $\mv{y}:=\mm{M}\mv{e}$ denotes a metric-weighted error vector, serving as a new optimization variable.
For discrete-time implementation, $\ddtt\mm{M}$ is approximated as $\tfrac{1}{T_s}(\mm{M}-\mm{M}_{\rm{pre}})$, where $T_s$ denotes the sampling time and $\mm{M}_{\rm{pre}}$ denotes the contraction metric obtained at the previous time step.

To justify the proposed approach, we first present the following theorem that provides a sufficient condition of the projected contraction condition \eqref{eq:cstr:ctrc:proj} to guarantee the full contraction condition \eqref{eq:cstr:ctrc}.

\begin{theorem} \label{thm:proj:ctrc:suff}
   For the system \eqref{eq:sys} satisfying Assumption~\ref{assum:sys:B} with the control law \eqref{eq:control:law}, the projected contraction condition \eqref{eq:cstr:ctrc:proj} is necessary for the full contraction condition \eqref{eq:cstr:ctrc}.
   Furthermore, suppose that there exists a metric-weighted error vector $\mv y=\mm M\mv e$ satisfying \eqref{eq:cstr:ctrc:proj}.
   If the reconstructed matrix $\mm M$ from $\mv y$ has a sufficiently large error-directional metric scale $\mu$ and a sufficiently small condition number $\chi$, then the reconstructed matrix also satisfies the full contraction condition \eqref{eq:cstr:ctrc}.
\end{theorem}

\begin{proof}

   \emph{Necessity.}

   If the full contraction condition \eqref{eq:cstr:ctrc} holds, then pre- and post-multiplying it by $\mv e^\top$ and $\mv e$ directly yields the projected condition \eqref{eq:cstr:ctrc:proj}.

   Therefore, the projected condition is necessary for the full contraction condition.

   \emph{Sufficiency.}

   Since the control effectiveness matrix $\mm{B}$ is full row rank matrix according to Assumption~\ref{assum:sys:B}, we have 
   \begin{equation}
      \mm{B}\mm{R}^{-1}\mm{B}^\top 
      \ge
      \sigma \mm{I}_n
      ,
   \end{equation}
   for some $\sigma\in\R_{>0}$.
   In addition, according to Assumption~\ref{assum:sys:smooth}, the SDC $\mathbb{A}$ is bounded such that $\norm{\mathbb{A}} \le \overline{A},$ for some $\overline{A}\in\R_{>0}$.

   The error-directional metric scale $\mu$ can be defined as follows:
   \begin{equation}
      \mu
      :=
      \frac{\mv{e}^\top\mm{M}\mv{e}}{\mv{e}^\top\mv{e}}
      .
   \end{equation}
   By the definition of $\mm{M}$, the following inequality holds:
   \begin{equation} \label{eq:mu:bounds}
      \underline{m}
      \le
      \mu
      \le
      \overline{m}
      \implies
      \tfrac{\mu}{\chi}
      \le 
      \underline{m}
      ,\ 
      \overline{m}
      \le
      \mu \chi
      ,\ 
      1
      \le 
      \chi 
      .
   \end{equation}
   % \mu/\underline{m} \le \chi
   % \mu / \chi \le \underline{m}
   % 1/\chi \le \mu/\overline{m}
   % \overline{m} \le \mu \chi

   Each term of the projected contraction condition \eqref{eq:cstr:ctrc:proj} can be conservatively upper-bounded as follows:
   \begin{align}
      -2
      \norm{
         \mv{y}^\top 
         \!
         \left(
         \!
            \mm{B} 
            \mm{R}^{-1} 
            \mm{B}^\top
         \!
         \right)
         \mv{y}
      }
      &\le
      -2\sigma\underline{m}^2
      \norm{\mv{e}}^2
      \\
      \norm{
         \mv{e}^\top
         \!
         \left(
            \tfrac{1}{T_s}+2\alpha
            +
            2\mathbb{A}^\top
         \!
         \right)
         \mv{y}
      }
      &\le
      \left(
         \tfrac{1}{T_s}
         +
         2
         \left(
            \alpha
            +
            \overline{A}
         \right)
      \right)
      \overline{m}
      \norm{\mv{e}}^2
      ,
      \\
      -
      \norm{
         \tfrac{\mv{e}^\top \mm{M}_{\rm{pre}} \mv{e}}{T_s}
      }
      &\le
      -
      \tfrac{1}{T_s}\underline{m}
      \norm{\mv{e}}^2
      . 
   \end{align}
   By substituting the above inequalities into \eqref{eq:cstr:ctrc:proj}, we have
   \begin{equation}
      \left(
      -2\sigma\underline{m}^2
      % \norm{\mv{e}}^2
      +
      \left(
         \tfrac{1}{T_s}
         +
         2
         \left(
            \alpha
            +
            \overline{A}
         \right)
      \right)
      \overline{m}
      -\tfrac{1}{T_s}\underline{m}
      \right) 
      \norm{\mv{e}}^2
      \le 0
      .
   \end{equation}
   Additionally, invoking \eqref{eq:mu:bounds}, the above inequality can be further rearranged as follows:
   \begin{equation}
      -2\sigma \tfrac{\mu^2}{\chi^2}
      +
      \left(
         \tfrac{1}{T_s}
         +
         2
         \left(
            \alpha
            +
            \overline{A}
         \right)
      \right)
      \mu \chi
      -\tfrac{1}{T_s}\tfrac{\mu}{\chi}
      \le 0
      .
   \end{equation}
   
   Finally we have the following sufficient condition for the full contraction condition \eqref{eq:cstr:ctrc}:
   \begin{align} \label{eq:suff:cond:mu}
      \mu\ge
      &
      \tfrac{1}{2\sigma}
      \left(
         2(\alpha+\overline{A})\chi^3
         +
         \tfrac{1}{T_s}
         \left(
            \chi^3
            -
            \chi
         \right)
      \right)
      .
   \end{align}
   The right-hand side of \eqref{eq:suff:cond:mu} is monotonically increasing with respect to $\chi$ for $\chi\ge 1$.
   Thus, if $\mu$ is sufficiently large and $\chi$ is sufficiently small such that the inequality \eqref{eq:suff:cond:mu} holds, the full contraction condition \eqref{eq:cstr:ctrc} is guaranteed.

\end{proof}

\subsection{Vector-Space Optimization Framework} \label{sec:sub:vector:space:opt}

According to Theorem~\ref{thm:proj:ctrc:suff}, the following optimization problem can be formulated, as follows:
\begin{subequations} \label{eq:opt:main}
   \begin{align}
      \min_{
         \mv{y},s
         } 
         \quad &
      \norm{\mv{y}}\norm{\mv{e}} - \mv{y}^\top \mv{e}
      % J(\mv{y})
      +
      \lambda s
      \label{eq:opt:main:obj}
      \\
      \text{subject to} 
      \quad &
      % \begin{cases}
         \eqref{eq:cstr:ctrc:proj}
         ,
         \label{eq:opt:main:ctrc:proj}
         \\ &
         \left(
            \mv{e}^\top \mv{y} 
            -
            \mu \mv{e}^\top \mv{e}
         \right)^2
         \le
         (\epsilon + s) \norm{\mv{e}}^2
         ,
         \label{eq:opt:main:energy}
         \\ &
         c(\mv{y}) \le 0
         ,
         \label{eq:opt:main:sat}
      % \end{cases}
   \end{align}
\end{subequations}
where $s,\epsilon\in\R_{\ge0}$ denote the slack variable and the allowable small energy variation, respectively of the energy-based constraint \eqref{eq:opt:main:energy} which will be explained later; $\lambda\in\R_{>0}$ denotes the penalty term for the slack variable $s$ and $c(\cdot)$ denotes input constraints, which can be designed to handle the input saturation effect.

The objective function \eqref{eq:opt:main:obj} is designed to minimize the angle between $\mv{y}$ and $\mv{e}$, which is equivalent to minimizing the condition number of the contraction metric $\mm{M}$.
Additionally, the last term in \eqref{eq:opt:main:obj} is a penalty term for the slack variable $s$ which will be explained later.

Imposed constraints include the projected contraction constraint \eqref{eq:opt:main:ctrc:proj}, the energy-based constraint \eqref{eq:opt:main:energy}, and the input saturation constraint \eqref{eq:opt:main:sat}.

The energy-based constraint \eqref{eq:opt:main:energy} is designed to prescribe the error-directional metric scale by enforcing 
$\mv e^\top\mv y\approx\mu\mv e^\top\mv e$; since $\mv y=\mm M\mv e$, this corresponds to 
$\mv e^\top\mm M\mv e/\mv e^\top\mv e\approx\mu$.
According to Theorem~\ref{thm:proj:ctrc:suff}, a sufficiently large error-directional metric scale, together with a sufficiently small condition number, promotes satisfaction of the full contraction condition \eqref{eq:cstr:ctrc}.
However, due to the input saturation constraint \eqref{eq:opt:main:sat}, the projected contraction constraint \eqref{eq:opt:main:ctrc:proj} may conflict with the input saturation constraint \eqref{eq:opt:main:sat}, leading to potential infeasibility of the optimization problem \eqref{eq:opt:main}.
To resolve this issue, the slack variable $s$ is introduced to relax the energy-based constraint.

The input saturation constraint \eqref{eq:opt:main:sat} can be designed to directly handle the input saturation effect by defining $c(\mv{y})$ to represent the admissible input set induced by the saturation function $\mysat(\cdot)$.

After the optimization, the contraction metric $\mm{M}$ can be retrieved while maintaining the condition number and the relationship $\mv{y}=\mm{M}\mv{e}$ as follows:
\begin{equation} \label{eq:retrive:mm}
   \mm{M}
   =
   \tfrac{1}{\sqrt{\chi}}
   \mydiag
   \left(
      \begin{pmatrix}
         \sqrt{\chi}, 1, \ldots, 1, \tfrac{1}{\sqrt{\chi}}
      \end{pmatrix}
   \right)
   \cdot
   \tfrac{1}{\norm{\mv{e}}\norm{\mv{y}}}
   ,
\end{equation}
where $\chi=\tfrac{1+\sin(\theta)}{1-\sin(\theta)}$ denotes the minimum condition number of $\mm{M}$ when the angle $\theta$ between $\mv{y}$ and $\mv{e}$ is given.

% In summary, the proposed formulation replaces the online search over a matrix-valued metric with the optimization of the feedback-relevant metric action $\mv y=\mm M\mv e$.
% The obtained vector $\mv y$ directly determines the constrained feedback input, while a positive definite matrix $\mm M$ satisfying $\mm M\mv e=\mv y$ is reconstructed for evaluating the projected condition and for use in the subsequent time step.
% Thus, the proposed framework should be interpreted as a projected contraction-based design that prioritizes scalable computation and direct input-constraint handling, while recovering the full contraction condition when the sufficient metric-scale and condition-number conditions in Theorem~\ref{thm:proj:ctrc:suff} are satisfied.

In conclusion, instead of solving the matrix-space optimization problem \eqref{eq:opt:conv:lmi}, the proposed method solves the vector-space optimization problem \eqref{eq:opt:main} to obtain the projected contraction metric $\mv{y}$.
After that, to obtain the contraction metric $\mm{M}$ for the control law \eqref{eq:control:law} and $\mm{M}_{\rm{pre}}$ in \eqref{eq:cstr:ctrc}, the contraction metric $\mm{M}$ is retrieved using \eqref{eq:retrive:mm}.
   
The overall procedure of the proposed method is summarized in Algorithm~\ref{alg:alg:1}, in which the optimization problem \eqref{eq:opt:main} is solved at each time step.

% \begin{algorithm}[!t]
%    \caption{Proposed controller} \label{alg:alg:1}
%    \begin{algorithmic}[1]
%       \State \textbf{Initialize:} $\mv{y}_0=\mv{1}_{n\times 1}$, $s_0=0$, $\mm{M}_{0}=\mu\mm{I}_n$
%       \For{$k = 1 $ to $\infty$}
%             \State Get current $\mv{x}_k$, $\mv{x}_{d,k}$, $\mv{u}_{d,k}$ and $\mv{e}_k$
%             \State Assign $\mv{y}_k \leftarrow \mv{y}_{k-1}$, $s_k \leftarrow s_{k-1}$ for warm-start
%             \State Solve \eqref{eq:opt:main} to get optimal $\mv{y}_k$ and $s_k$
%             \State Calculate $\mm{M}_{k}$ using $\mv{y}_k$ and $\mv{e}_k$ using \eqref{eq:retrive:mm}
%             \State Apply the control law \eqref{eq:control:law} using $\mm{M}_{k}$
%       \EndFor
%    \end{algorithmic}
% \end{algorithm}

\subsection{Illustrative Example in Two-Dimensional Space} \label{sec:sub:illustrative}

% \begin{figure}[t]
%    \centering
%    \includegraphics[width=0.99\linewidth]
%    {
%       figures/scene_first.drawio.pdf
%    }
%    \caption{
%       Illustration of the optimization problem \eqref{eq:opt:main} in two-dimensional space; see Section \ref{sec:sub:illustrative} for details.
%    }
% \label{fig:scene:first}
% \end{figure}

To facilitate geometric interpretation, the optimization problem \eqref{eq:opt:main} is illustrated in two-dimensional space in Fig.~\ref{fig:scene:first}.
The projected contraction condition \eqref{eq:opt:main:ctrc:proj} and the input saturation constraint \eqref{eq:opt:main:sat} restrict the feasible space of $\mv{y}$ to the outside of the blue line and the inside of the red line, respectively.
The green line represents the energy-based constraint \eqref{eq:opt:main:energy}, assuming that the allowable energy variation $\epsilon$ is approximately zero.

% Since the projected contraction constraint \eqref{eq:opt:main:ctrc:proj} is an ellipse-shaped concave function and the input saturation constraint \eqref{eq:opt:main:sat} is a convex function, there exist at most two intersection points of these constraints.
Without the input saturation constraint, the optimal solution is obtained at point (a) where \eqref{eq:opt:main:ctrc:proj} and \eqref{eq:opt:main:energy} are active, resulting in the minimum condition number of $\mm{M}$ and the error-directional metric scale $\mu$.
However, when the input saturation constraint is considered, the optimal solution is obtained at point (b) where \eqref{eq:opt:main:sat} and \eqref{eq:opt:main:energy} are active, resulting in relatively large condition number of $\mm{M}$ and insufficient error-directional metric scale $\mu$.
In this case, the slack variable $s$ is increased to relax the energy-based constraint \eqref{eq:opt:main:energy} to satisfy the input saturation constraint, leading to potential violation of sufficiency conditions in Theorem~\ref{thm:proj:ctrc:suff}.


\section{Validations}

\section{Conclusion}
\label{sec:conclusion}

Will be added later.

% ============================================
%                   Appendix
% ============================================
%% The Appendices part is started with the command \appendix;
%% appendix sections are then done as normal sections
\appendix

% ============================================
%                   Bibliography
% ============================================
\bibliographystyle{IEEEtran}
\bibliography{\template/refs}

% ============================================
%                   BIOGRAPHIES
% ============================================
\begin{IEEEbiography}[{
  \includegraphics[width=1in,height=1.25in,clip,keepaspectratio]{\template/people/msRyu.jpg}}]{Myeongseok Ryu}
  received the B.S. degree in mechanical engineering from Incheon National University, South Korea, in 2023, and the M.S. degree in mechanical engineering from Gwangju Institute of Science and Technology (GIST), South Korea, in 2025. 
  He is currently serving as a researcher at the Cho Chun Shik Graduate School of Mobility, Korea Advanced Institute of Science and Technology (KAIST), South Korea.
  His research interests include adaptive control, neural-networks, and constrained optimization.
\end{IEEEbiography}

\begin{IEEEbiography}[{
  \includegraphics[width=1in,height=1.25in,clip,keepaspectratio]{\template/people/khChoi.jpg}}]{Kyunghwan Choi}
  Kyunghwan Choi received his B.S., M.S., and Ph.D. degrees in Mechanical Engineering from KAIST, Daejeon, South Korea, in 2014, 2016, and 2020, respectively. Following his Ph.D., he worked at the Center for Eco-Friendly \& Smart Vehicles at KAIST as a Postdoctoral Fellow and later as a Research Assistant Professor. In 2022, he joined GIST as an Assistant Professor in the Department of Mechanical and Robotics Engineering. Currently, he serves as an Assistant Professor at the Cho Chun Shik Graduate School of Mobility at KAIST, where he is also the Director of the Mobility Intelligence and Control Laboratory. His research focuses on optimal and learning-based control for connected, automated, and electrified vehicles (CAEVs).
\end{IEEEbiography}

\end{document}